% Source: physical PDF pp. 349--356 (printed pp. 326--333), from the
% Section 31.2 heading through its closing sentence immediately before
% Section 31.3.
% Inventory: Eqs. (31.2.1)--(31.2.34), eleven unnumbered displays, six
% linked citation markers (references [6], [7], [7a], [8], [9], and [10]),
% one centered three-asterisk divider, and no figures, tables, or footnotes.
% Coverage-manifest discrepancy: the source has a centered three-asterisk
% divider on physical PDF p. 354, although COVERAGE.md does not list one
% in Section 31.2.
% Eq. (31.2.16) retains the source's unsimplified coefficients 6/8 and 6/16.
% Corrected the measure and missing integral in Eq. (31.2.27).
% Uncertainties: none.

\subsection{The Gravitational Action}

To find a suitable gravitational action, we must construct a superfield that
is invariant under the generalized gauge transformation
\(H_\mu\to H_\mu+\Delta_\mu\). As a starting point, we recall that the field
\(b_\mu\) defined by Eq.~\eqref{eq:31.1.27} is invariant under this
transformation. By making successive supersymmetry transformations we can see
that \(b_\mu\) is the \(C\)-component of an `Einstein' superfield \(E_\mu\),
whose components are
\begin{equation}
  C_\mu^E=b_\mu\,,
  \tag{31.2.1}\label{eq:31.2.1}
\end{equation}
\begin{equation}
  \omega_{\mu\alpha}^E
  =\frac32L_{\mu\alpha}
   +\frac12(\sigma_\mu\bar\sigma^\nu L_\nu)_\alpha\,,
  \qquad\text{and its complex conjugate},
  \tag{31.2.2}\label{eq:31.2.2}
\end{equation}
\begin{equation}
  M_\mu^E=\partial_\mu s,\qquad
  N_\mu^E=\partial_\mu p\,,
  \tag{31.2.3}\label{eq:31.2.3}
\end{equation}
\begin{equation}
  V_{\mu\nu}^E
  =-\frac32E_{\mu\nu}
   +\frac12\eta_{\mu\nu}E^\rho{}_\rho
   +\frac12\epsilon_{\nu\mu\sigma\rho}\partial^\sigma b^\rho\,,
  \tag{31.2.4}\label{eq:31.2.4}
\end{equation}
\begin{equation}
  \lambda_{\mu\alpha}^E
  =-\ii\partial_\mu
     (\sigma^\nu\bar\omega_\nu^E)_\alpha
   +\ii(\sigma^\rho\partial_\rho\bar\omega_\mu^E)_\alpha\,,
  \qquad\text{and its complex conjugate},
  \tag{31.2.5}\label{eq:31.2.5}
\end{equation}
\begin{equation}
  D_\mu^E=\partial_\mu\partial^\nu b_\nu-\Box b_\mu\,,
  \tag{31.2.6}\label{eq:31.2.6}
\end{equation}
where \(E_{\mu\nu}\) is the linearized Einstein tensor
\begin{equation}
  \begin{aligned}
  E_{\mu\nu}\equiv\frac12\bigl(
    &\partial_\mu\partial_\nu h^\lambda{}_\lambda
     +\Box h_{\mu\nu}
     -\partial_\mu\partial^\lambda h_{\lambda\nu}
     -\partial_\nu\partial^\lambda h_{\lambda\mu}
  \\
    &-\eta_{\mu\nu}\Box h^\lambda{}_\lambda
     +\eta_{\mu\nu}\partial^\lambda\partial^\rho h_{\lambda\rho}
  \bigr)
  \end{aligned}
  \tag{31.2.7}\label{eq:31.2.7}
\end{equation}
and
\begin{equation}
  L^\nu_\alpha\equiv
  \epsilon^{\nu\mu\kappa\rho}
  \sigma_{\mu\,\alpha\dot\alpha}
  \partial_\kappa\bar\psi_\rho^{\dot\alpha}\,,
  \qquad
  \bar L^{\nu\dot\alpha}\equiv(L^\nu_\alpha)^*,
  \tag{31.2.8}\label{eq:31.2.8}
\end{equation}
which will be shown in Section 31.3 to be the left-hand side of the wave
equation for a free massless field of spin \(3/2\). For instance, by applying
the supersymmetry transformation rules \eqref{eq:26.2.11},
\eqref{eq:26.2.15}, and \eqref{eq:26.2.17} for \(C^H\), \(V_\mu^H\), and
\(D^H\) to Eqs.~\eqref{eq:31.2.1} and \eqref{eq:31.1.27} we find
\[
  \begin{aligned}
  \delta C^{E\sigma}
  =\ii\bigl(\alpha^\beta K^\sigma_\beta
  -\bar\alpha_{\dot\beta}\bar K^{\sigma\dot\beta}\bigr),\qquad
  K^\sigma_\alpha={}&
    -\ii(\sigma^\rho\partial_\rho\bar\lambda^{H\sigma})_\alpha\\
  &-\frac12\epsilon^{\nu\mu\kappa\sigma}
    \sigma_{\mu\,\alpha\dot\alpha}
    \partial_\kappa\bar\lambda_\nu^{H\dot\alpha}
    +\partial^\sigma\partial^\kappa\omega_{\kappa\alpha}^H ,
  \end{aligned}
\]
where \(\bar K^{\sigma\dot\alpha}=(K^\sigma_\alpha)^*\).
Comparison of this with the transformation rule \eqref{eq:26.2.11} for
\(C^E\) shows that
\[
  \omega^{E\sigma}_\alpha
  =-\ii(\sigma^\rho\partial_\rho\bar\lambda^{H\sigma})_\alpha
   -\frac12\epsilon^{\nu\mu\kappa\sigma}
    \sigma_{\mu\,\alpha\dot\alpha}
    \partial_\kappa\bar\lambda_\nu^{H\dot\alpha}
   +\partial^\sigma\partial^\kappa\omega_{\kappa\alpha}^H\,,
\]
We can use Eq.~\eqref{eq:31.1.26} to express \(\lambda_\mu^H\) in terms of
\(\psi_\mu\) and \(\omega_\mu^H\):
\[
  \lambda_{\mu\alpha}^H
  =\psi_{\mu\alpha}
   +(\sigma_\mu\bar\sigma^\rho\psi_\rho)_\alpha
   +\ii\sigma_{\mu\,\alpha\dot\alpha}
     \partial^\rho\bar\omega_\rho^{H\dot\alpha}\,,
\]
and find that \(\omega_\sigma^E\) may be written in terms of \(\psi_\sigma\)
alone
\[
  \omega_{\sigma\alpha}^E
  =-\ii(\sigma^\rho\partial_\rho\bar\psi_\sigma)_\alpha
   +\ii\partial_\sigma
      (\sigma^\rho\bar\psi_\rho)_\alpha
   -\frac12\epsilon_{\nu\mu\kappa\sigma}
    \sigma^{\mu}{}_{\alpha\dot\alpha}
    \partial^\kappa\bar\psi^{\nu\dot\alpha}\,.
\]
Using the identity
\[
  -\ii\bigl(\eta_{\mu\nu}\sigma_\lambda
          -\eta_{\mu\lambda}\sigma_\nu\bigr)
  =\epsilon_{\mu\lambda\nu\rho}\sigma^\rho
   +\frac12\epsilon_{\sigma\lambda\nu\rho}
     \sigma_\mu\bar\sigma^\sigma\sigma^\rho\,,
\]
we find Eq.~\eqref{eq:31.2.2} for \(\omega_\mu^E\). Continuing in the same
way gives the other formulas \eqref{eq:31.2.3}--\eqref{eq:31.2.6} for the
components of \(E_\mu\) and confirms that these components make up a real
superfield.

We can now form a quadratic action for the metric superfield that is invariant
under both supersymmetry and the extended gauge transformation
\(H_\mu\to H_\mu+\Delta_\mu\), by taking the Lagrangian density as
\begin{equation}
  \mathcal L_E
  \equiv\frac43(E_\mu H^\mu)_D
  =E_{\mu\nu}h^{\mu\nu}
   -\frac12\bigl(\psi_\mu L^\mu
                 +\bar\psi_\mu\bar L^\mu\bigr)
   -\frac43(s^2+p^2-b_\mu b^\mu)\,.
  \tag{31.2.9}\label{eq:31.2.9}
\end{equation}
The factor \(4/3\) is chosen to give the kinematic Lagrangian for the graviton
field a conventional sign and normalization: apart from terms involving
\(\partial^\mu h_{\mu\nu}\) or \(h^\lambda{}_\lambda\), it is a sum of
Klein--Gordon Lagrangians for the components of \(h_{\mu\nu}\). The
normalization of the gravitino field is discussed in the next section.

To see that the first two terms in the final expression in
Eq.~\eqref{eq:31.2.9} are invariant under the gauge transformations
\eqref{eq:31.1.12} and \eqref{eq:31.1.13}, we note that \(E_{\mu\nu}\) and
\(L_\mu\) are manifestly invariant under these transformations, and that the
action is symmetric between the two factors of \(h_{\mu\nu}\) or \(\psi_\mu\)
appearing in these terms. The absence in Eq.~\eqref{eq:31.2.9} of derivatives
of \(s\), \(p\), and \(b_\mu\) shows that these are auxiliary fields. The field
equations make them vanish for pure gravitation, but not when gravitation is
coupled to matter.

Before taking up the coupling of matter and gravitation, let us pause to
consider what value we should give the normalization constant \(\kappa\) in
Eqs.~\eqref{eq:31.1.5} and \eqref{eq:31.1.10}. The Einstein--Hilbert action
for a pure gravitational field is
\begin{equation}
  I_{\mathrm{GR}}
  =-\frac{1}{16\pi G}\int \dd^4x\,\sqrt{g}\,R\,,
  \tag{31.2.10}\label{eq:31.2.10}
\end{equation}
where \(G\) is the Newton constant of gravitation, \(g(x)\) is the determinant
of the metric tensor \(g_{\mu\nu}(x)\), and \(R(x)\) is the curvature scalar
calculated from \(g_{\mu\nu}(x)\). To calculate \(I_{\mathrm{GR}}\) for a weak
gravitational field with
\(g_{\mu\nu}=\eta_{\mu\nu}+2\kappa h_{\mu\nu}\), we may recall that for
arbitrary variations \(\delta g_{\mu\nu}(x)\) in the gravitational
field~\hyperref[ch31-ref-6]{[6]}
\begin{equation}
  \delta I_{\mathrm{GR}}
  =\frac{1}{16\pi G}\int \dd^4x\,\sqrt{g}\,
   \left[R^{\mu\nu}-\frac12g^{\mu\nu}R\right]\delta g_{\mu\nu}\,,
  \tag{31.2.11}\label{eq:31.2.11}
\end{equation}
where \(R^{\mu\nu}(x)\) is the Ricci tensor calculated from
\(g_{\mu\nu}(x)\). For a weak field with
\(g_{\mu\nu}=\eta_{\mu\nu}+2\kappa h_{\mu\nu}\), the Ricci tensor
is~\hyperref[ch31-ref-7]{[7]}
\begin{equation}
  R^{\mu\nu}
  =\kappa\left(
     \Box h^{\mu\nu}
     -\partial_\lambda\partial^\mu h^{\lambda\nu}
     -\partial_\lambda\partial^\nu h^{\lambda\mu}
     +\partial^\mu\partial^\nu h^\lambda{}_\lambda
   \right),
  \tag{31.2.12}\label{eq:31.2.12}
\end{equation}
so for weak fields
\begin{equation}
  R^{\mu\nu}-\frac12g^{\mu\nu}R=2\kappa E^{\mu\nu}\,,
  \tag{31.2.13}\label{eq:31.2.13}
\end{equation}
and therefore Eq.~\eqref{eq:31.2.10} gives
\[
  \delta I_{\mathrm{GR}}
  =\frac{\kappa^2}{4\pi G}
   \int \dd^4x\,E^{\mu\nu}\delta h_{\mu\nu}\,.
\]
On the other hand, taking account of the symmetry of
\(\int \dd^4x\,E^{\mu\nu}h_{\mu\nu}\) between the two factors of \(h\) it
contains, we have
\[
  \delta\int \dd^4x\,E^{\mu\nu}h_{\mu\nu}
  =2\int \dd^4x\,E^{\mu\nu}\delta h_{\mu\nu}\,.
\]
In order for the term \(E_{\mu\nu}h^{\mu\nu}\) in
Eq.~\eqref{eq:31.2.9} to give the usual gravitational Lagrangian density, it
is therefore necessary to take
\begin{equation}
  \kappa=\sqrt{8\pi G}\,.
  \tag{31.2.14}\label{eq:31.2.14}
\end{equation}

Let's now combine the Einstein Lagrangian density \eqref{eq:31.2.9} with the
interaction\newline
\eqref{eq:31.1.34} between gravitation and matter and with the
matter Lagrangian \(\mathcal L_M\), to form a total Lagrangian density:
\begin{equation}
  \begin{aligned}
  \mathcal L={}&
    \mathcal L_M+E_{\mu\nu}h^{\mu\nu}
    -\frac12\bigl(\psi_\mu L^\mu
                  +\bar\psi_\mu\bar L^\mu\bigr)
    -\frac43(s^2+p^2-b_\mu b^\mu)
  \\
  &+2\kappa\left[
      \frac12\mathcal R_\sigma b^\sigma
      +\frac14\bigl(S^\sigma\psi_\sigma
                    +\bar S^\sigma\bar\psi_\sigma\bigr)
      -\mathcal M s-\mathcal Np
      +\frac12T^{\kappa\sigma}h_{\sigma\kappa}
    \right].
  \end{aligned}
  \tag{31.2.15}\label{eq:31.2.15}
\end{equation}
The field equations for the auxiliary fields give
\begin{equation}
  s=-6\kappa\mathcal M/8,\qquad
  p=-6\kappa\mathcal N/8,\qquad
  b_\mu=-6\kappa\mathcal R_\mu/16\,.
  \tag{31.2.16}\label{eq:31.2.16}
\end{equation}
Using this to eliminate the auxiliary fields, the Lagrangian density
\eqref{eq:31.2.15} now gives
\begin{equation}
  \begin{aligned}
  \mathcal L={}&
    \mathcal L_M+E_{\mu\nu}h^{\mu\nu}
    -\frac12\bigl(\psi_\mu L^\mu
                  +\bar\psi_\mu\bar L^\mu\bigr)
    +\frac34\kappa^2
     \left(\mathcal M^2+\mathcal N^2
           -\frac14\mathcal R_\mu\mathcal R^\mu\right)
  \\
  &+\frac12\kappa
    \bigl(S^\sigma\psi_\sigma+\bar S^\sigma\bar\psi_\sigma\bigr)
    +\kappa T^{\kappa\sigma}h_{\sigma\kappa}\,.
  \end{aligned}
  \tag{31.2.17}\label{eq:31.2.17}
\end{equation}
The sources of the fields \(\psi_\mu\) and \(h_{\mu\nu}\) are of order
\(\kappa\), so we can regard these fields as being of this order, which makes
all terms in Eq.~\eqref{eq:31.2.17} beyond \(\mathcal L_M\) of order
\(\kappa^2\).

In the vacuum only the scalar fields \(s\) and \(p\) can have
tree-approximation expectation values, giving the vacuum an energy density,
to order \(G\):
\begin{equation}
  \rho_{\mathrm{VAC}}
  =-\mathcal L_{\mathrm{VAC}}
  =-\mathcal L_{M\,\mathrm{VAC}}
   -\frac34\kappa^2(\mathcal M^2+\mathcal N^2)\,.
  \tag{31.2.18}\label{eq:31.2.18}
\end{equation}
The negative sign of the term in the vacuum energy density of first order in
\(G\) is a characteristic difference between theories of supergravity and
ordinary supersymmetry.

For instance, in the theory of a single left-chiral superfield \(\Phi\) with
superpotential \(f(\Phi)\), the zeroth-order vacuum energy is
\(-\mathcal L_{M\,\mathrm{VAC}}=\lvert \dd f(\phi)/\dd\phi\rvert^2\), while
Eqs.~\eqref{eq:31.1.33} and \eqref{eq:26.7.27} give
\[
  \mathcal M+\ii\mathcal N
  =-\frac13\left[
    \phi\frac{\dd f(\phi)}{\dd\phi}-3f(\phi)
  \right],
\]
so to first order in \(G\), the total vacuum energy is
\begin{equation}
  \rho_{\mathrm{VAC}}
  =\left\lvert\frac{\dd f(\phi)}{\dd\phi}\right\rvert^2
   -\frac{8\pi G}{3}
    \left\lvert
      \phi\frac{\dd f(\phi)}{\dd\phi}-3f(\phi)
    \right\rvert^2,
  \tag{31.2.19}\label{eq:31.2.19}
\end{equation}
where \(\phi\) is the scalar component of \(\Phi\). This is for the definition
of the metric for which the energy-momentum tensor is given by
Eq.~\eqref{eq:26.7.42}, so that in particular \(T^\lambda{}_\lambda\)
vanishes for \(f=0\). For other definitions the vacuum energy would be
changed by terms of order
\(8\pi G\lvert\phi\rvert^2\lvert \dd f(\phi)/\dd\phi\rvert^2\). However, this
ambiguity is unimportant in calculating the minimum value of the vacuum
energy. Inspection of Eq.~\eqref{eq:31.2.19} shows that if \(f(\phi)\) is
stationary at some point \(\phi_0\), then \(\rho_{\mathrm{VAC}}\) has a local
minimum for \(\lvert\phi-\phi_0\rvert\) and \(\lvert \dd f/\dd\phi\rvert\) of first
order in \(G\), and at any such point the vacuum energy to first order in
\(G\) is
\begin{equation}
  \rho_{\mathrm{VAC}}=-24\pi G\lvert f(\phi_0)\rvert^2\,.
  \tag{31.2.20}\label{eq:31.2.20}
\end{equation}
There is an algebraic reason why vacuum states with unbroken supersymmetry in
supergravity theories cannot have positive energy density.

The solutions of the Einstein field equations for a uniform non-zero vacuum
energy density \(\rho_V\) take the form of a de Sitter space for
\(\rho_V>0\) and anti-de Sitter space for \(\rho_V<0\). These spaces may be
described as the surfaces
\begin{equation}
  x_5^2\mathbin{\pm}\eta_{\mu\nu}x^\mu x^\nu=R^2\,,
  \tag{31.2.21}\label{eq:31.2.21}
\end{equation}
in a quasi-Euclidean five-dimensional space with line element
\begin{equation}
  \dd s^2=\eta_{\mu\nu}\dd x^\mu \dd x^\nu\mathbin{\pm}\dd x_5^2\,,
  \tag{31.2.22}\label{eq:31.2.22}
\end{equation}
with the upper sign for de Sitter space and the lower sign for anti-de Sitter
space. The spacetime symmetry of these spaces is no longer the Poincaré group
consisting of translations and Lorentz transformations, but instead \(O(4,1)\)
for de Sitter space and \(O(3,2)\) for anti-de Sitter space, where \(O(n,m)\)
is the group of linear transformations that leave invariant a diagonal metric
with \(n\) positive and \(m\) negative elements on the main diagonal. The
supersymmetry that is unbroken in theories with a de Sitter or anti-de Sitter
vacuum state therefore has the simple group \(O(4,1)\) or \(O(3,2)\),
respectively, as its spacetime symmetry. Nahm~\hyperref[ch31-ref-7a]{[7a]} has
cataloged all supersymmetries with simple spacetime symmetries. There are
simple \(O(3,2)\) supersymmetries, as well as \(N\)-extended \(O(3,2)\)
supersymmetries, but for \(O(4,1)\) there is only an \(N=2\) supersymmetry. We
have been considering \(N=1\) supergravity theories here, so they can have
vacuum states with unbroken supersymmetry and \(\rho_V<0\), which gives
\(O(3,2)\) spacetime symmetry, but not with unbroken supersymmetry and
\(\rho_V>0\), which would give \(O(4,1)\) spacetime symmetry.

The possibility of vacuum field configurations with negative energy density
may seem at first sight to threaten the stability of our universe. It is
common for \(f(\phi)\) to have several stationary points, at which it takes
different values. Even if we fine-tune the parameters in \(f(\phi)\) so that
\(f(\phi)\) vanishes at one of these stationary points, accounting for the
observed nearly flat space of our universe, any other stationary point with a
non-zero value of \(f(\phi)\) will yield a lower vacuum energy, raising the
possibility of a collapse to a state of negative energy density, with a metric
that would have to be of the `anti-de Sitter' form rather than flat.

Fortunately, the value \eqref{eq:31.2.20} is just barely insufficient to
compensate for the positive energy that has to go into surface tension in
making a bubble of anti-de Sitter space, so that the transition from ordinary
flat space is not actually favored energetically. Coleman and de
Luccia~\hyperref[ch31-ref-8]{[8]} have applied the equations of general
relativity to a bubble of negative internal energy density \(-\epsilon\) and
positive surface tension \(S\) in a flat space of energy density zero, and
have shown that any such bubble that does not entail gravitational
singularities will have positive energy if
\begin{equation}
  \epsilon\leq 6\pi GS^2\,.
  \tag{31.2.23}\label{eq:31.2.23}
\end{equation}
The surface tension is the energy per area in the bubble surface, given to
zeroth order in \(G\) by the integral of the energy density through the bubble
wall:
\begin{equation}
  S_1
  =\int_{r_-}^{r_+}\dd r
   \left[
     \left\lvert\frac{\dd\phi}{\dd r}\right\rvert^2
     +\left\lvert\frac{\dd f(\phi)}{\dd\phi}\right\rvert^2
   \right],
  \tag{31.2.24}\label{eq:31.2.24}
\end{equation}
with \(r_-\) and \(r_+\) taken just inside and just outside the bubble wall.
This can be rewritten as
\[
  S_1
  =\int_{r_-}^{r_+}\dd r
   \left[
     \left\lvert
       \frac{\dd\phi}{\dd r}
       +\xi\left(\frac{\dd f(\phi)}{\dd\phi}\right)^*
     \right\rvert^2
     -2\operatorname{Re}\left(
       \xi^*\frac{\dd\phi}{\dd r}\frac{\dd f(\phi)}{\dd\phi}
     \right)
   \right],
\]
where \(\xi\) is an arbitrary phase factor with \(\lvert\xi\rvert=1\). The
integral in the second term is trivial: since we assume that \(\phi(r_+)\) is
at a value where \(f(\phi)\) is stationary and vanishes, while \(\phi(r_-)\)
is at some value \(\phi_0\) where \(f(\phi)\) is stationary but non-zero, the
integral is
\[
  \int_{r_-}^{r_+}\dd r\,
  \frac{\dd\phi}{\dd r}\frac{\dd f(\phi)}{\dd\phi}
  =-f(\phi_0)\,.
\]
To maximize this term, we take
\(\xi=f(\phi_0)/\lvert f(\phi_0)\rvert\), and obtain the
inequality~\hyperref[ch31-ref-9]{[9]}
\begin{equation}
  S_1\geq 2\lvert f(\phi_0)\rvert\,,
  \tag{31.2.25}\label{eq:31.2.25}
\end{equation}
with equality attained only if (as is usual) there is a solution of the
differential equation
\(\dd\phi/\dd r=-\xi(\dd f/\dd\phi)^*\) with the appropriate boundary conditions. The
inequality \eqref{eq:31.2.23} is thus satisfied if the internal energy density
is not less than \(-24\pi G\lvert f(\phi_0)\rvert^2\), which is precisely the
value \eqref{eq:31.2.20}. This calculation leaves open the possibility of an
instability of flat space due to radiative corrections, but the reader need
not worry: it has been proved that in supergravity theories in which there is
a vacuum field configuration with zero vacuum energy, the energy of any
disturbance in the fields which is limited to a finite region is
positive~\hyperref[ch31-ref-10]{[10]}.

\begin{center}
\(\ast\qquad\ast\qquad\ast\)
\end{center}

Eq.~\eqref{eq:26.7.48} shows that the energy-momentum tensor for a set of
left-chiral scalar superfields \(\Phi_n\) contains a term
\begin{equation}
  \Delta T^{\mu\nu}
  =\frac13
   \left(\eta^{\mu\nu}\Box-\partial^\mu\partial^\nu\right)
   \sum_n\lvert\phi_n\rvert^2\,.
  \tag{31.2.26}\label{eq:31.2.26}
\end{equation}
Integrating by parts, the corresponding interaction
\(\kappa h_{\mu\nu}\Delta T^{\mu\nu}\) makes a contribution to the action of
the form
\begin{equation}
  \frac{\kappa}{3}\int \dd^4x\,
  \sum_n\lvert\phi_n\rvert^2
  \left(\eta^{\mu\nu}\Box-\partial^\mu\partial^\nu\right)h_{\mu\nu}
  =\frac16\int \dd^4x\,R^{(1)}\sum_n\lvert\phi_n\rvert^2\,,
  \tag{31.2.27}\label{eq:31.2.27}
\end{equation}
where \(R^{(1)}\) is the curvature scalar in linear approximation. This is
added to the usual Einstein--Hilbert action
\(-\sqrt{g}R/2\kappa^2\) (which appears in Eq.~\eqref{eq:31.2.17} as the term
\(E_{\mu\nu}h^{\mu\nu}\)) and has the effect of replacing the coefficient of
this term with
\begin{equation}
  -\frac{1}{2\kappa^2}
  +\frac16\sum_n\lvert\phi_n\rvert^2
  =-\frac{1}{2\kappa^2}
   \left(1-\frac{\kappa^2}{3}\sum_n\lvert\phi_n\rvert^2\right).
  \tag{31.2.28}\label{eq:31.2.28}
\end{equation}
In order to restore the usual constant of gravitation, we can subject the
metric to a \emph{Weyl transformation}, replacing the vierbein \(e^a{}_\mu\)
with
\begin{equation}
  \widetilde e^a{}_\mu
  =e^a{}_\mu
   \sqrt{1-\frac{\kappa^2}{3}\sum_n\lvert\phi_n\rvert^2}\,.
  \tag{31.2.29}\label{eq:31.2.29}
\end{equation}
That is, we replace the metric \(g_{\mu\nu}\) with
\begin{equation}
  \widetilde g_{\mu\nu}
  =\left(1-\frac{\kappa^2}{3}\sum_n\lvert\phi_n\rvert^2\right)g_{\mu\nu}\,,
  \tag{31.2.30}\label{eq:31.2.30}
\end{equation}
or for weak fields
\begin{equation}
  \widetilde h_{\mu\nu}
  =h_{\mu\nu}
   -\frac{\kappa}{6}\sum_n\lvert\phi_n\rvert^2\eta_{\mu\nu}\,.
  \tag{31.2.31}\label{eq:31.2.31}
\end{equation}
The weak field Einstein tensor \eqref{eq:31.2.7} for the new metric is
\begin{equation}
  \begin{aligned}
  \widetilde E_{\mu\nu}\equiv\frac12\bigl(
    &\partial_\mu\partial_\nu\widetilde h^\lambda{}_\lambda
     +\Box\widetilde h_{\mu\nu}
     -\partial_\mu\partial^\lambda\widetilde h_{\lambda\nu}
     -\partial_\nu\partial^\lambda\widetilde h_{\lambda\mu}
  \\
    &-\eta_{\mu\nu}\Box\widetilde h^\lambda{}_\lambda
     +\eta_{\mu\nu}\partial^\lambda\partial^\rho
      \widetilde h_{\lambda\rho}
  \bigr)
  \\
  ={}&E_{\mu\nu}
      -\frac{\kappa}{6}
       \left(\partial_\mu\partial_\nu-\eta_{\mu\nu}\Box\right)
       \sum_n\lvert\phi_n\rvert^2\,.
  \end{aligned}
  \tag{31.2.32}\label{eq:31.2.32}
\end{equation}
The sum of the Einstein term in the original action and the term
\eqref{eq:31.2.27} is then
\begin{equation}
  \begin{aligned}
  &\int \dd^4x
   \left[
     h_{\mu\nu}E^{\mu\nu}
     +\frac{\kappa}{3}\sum_n\lvert\phi_n\rvert^2
      \left(\eta^{\mu\nu}\Box-\partial^\mu\partial^\nu\right)h_{\mu\nu}
   \right]
  \\
  &\qquad
  =\int \dd^4x
   \left[
     \widetilde h_{\mu\nu}\widetilde E^{\mu\nu}
     +\frac{\kappa^2}{12}
      \left(\partial_\mu\sum_n\lvert\phi_n\rvert^2\right)
      \left(\partial^\mu\sum_n\lvert\phi_n\rvert^2\right)
   \right],
  \end{aligned}
  \tag{31.2.33}\label{eq:31.2.33}
\end{equation}
so the effective gravitational constant is now actually constant. This
redefinition of the metric also produces a change in the potential. The
original Lagrangian density contains a term
\(-e\sum_n\lvert\partial f(\phi)/\partial\phi_n\rvert^2\), which in terms of
the new vierbein reads
\[
  \begin{aligned}
  -e\sum_n\left\lvert\frac{\partial f(\phi)}{\partial\phi_n}\right\rvert^2
  ={}&
  -\widetilde e
   \sum_n\left\lvert\frac{\partial f(\phi)}{\partial\phi_n}\right\rvert^2
  \\
  &-\frac{2\kappa^2}{3}
   \left(\sum_n\lvert\phi_n\rvert^2\right)
   \sum_n\left\lvert\frac{\partial f(\phi)}{\partial\phi_n}\right\rvert^2.
  \end{aligned}
\]
With the new definition of the metric, the potential \eqref{eq:31.2.19} is
therefore replaced with
\begin{equation}
  \begin{aligned}
  \rho_{\mathrm{VAC}}={}&
    \sum_n
    \left\lvert\frac{\partial f(\phi)}{\partial\phi_n}\right\rvert^2
    -\frac{\kappa^2}{3}
     \left\lvert
       \sum_n\phi_n\frac{\partial f(\phi)}{\partial\phi_n}
       -3f(\phi)
     \right\rvert^2
  \\
  &+\frac{2\kappa^2}{3}
    \left(\sum_n\lvert\phi_n\rvert^2\right)
    \sum_n
    \left\lvert\frac{\partial f(\phi)}{\partial\phi_n}\right\rvert^2.
  \end{aligned}
  \tag{31.2.34}\label{eq:31.2.34}
\end{equation}
The new term does not change the value of \(\rho_{\mathrm{VAC}}\) at its
stationary point to order \(\kappa^2\), so no change is needed in our previous
discussion of vacuum stability.
